% page 471 of the facsimile (image p-470 of the scan)
% block 172.7 x 259.3 mm ; left margin 29.3 mm
\pgc{7.620mm}{0.000mm}{
\l{}
\l{}
\l{}
\l{}
\l{\cel{56}463}
\l{}
\l{\cel{4}\sou{profilaxio}. (\sou{medicino)}. Rejimo qua prezervas kontre la malaji,}
\l{\cel{7}kontre la morbi. - DEFIRS.}
\l{}
\l{\cel{4}\sou{profitar}. (\sou{trans.}) I. Ektirar avantajo de ulo. - II. Obtenar}
\l{\cel{7}gano de ulo. - DEFIR.}
\l{}
\l{\cel{4}\sou{profunda}. I. Di qua la fundo esas tre infra kompare kun la ori-}
\l{\cel{7}fico, kun la surfaco. - II. Qua su extensas tre adfore, de lo}
\l{\cel{7}avana til lo dopa. - III. (\sou{metaf.}) (1) Qua exploras absolute}
\l{\cel{7}omna faktori di la kozi. (2) Qua esas extrema. - EFIS.}
\l{}
\l{\cel{4}\sou{prognoza}r.(\sou{trans}.) (\sou{medicino, meteorologio}). Konjektar, de ula}
\l{\cel{7}signi, to quo mustas eventor. - DEFIRS.}
\l{}
\l{\cel{4}\sou{programo}. I. Skriburo od imprimuro qua anuncas la detali di}
\l{\cel{7}ceremonio, di festo publika, di koncerto, du reprezento tea-}
\l{\cel{7}trala, e c. - II. Anunco di la temi quin profesoro traktos}
\l{\cel{7}en sua kurso. - III. Expozo di la koncepti politikala di par-}
\l{\cel{7}tiso, di individuo. - DEFIRS.}
\l{}
\l{\cel{4}\sou{progresar}. (\sou{netrans.,}\cel{1}ad). Avancar a grado superiora. - DEFIRS.}
\l{}
\l{\cel{4}\sou{progresiono}. (\sou{matem.)}\cel{2}Intersequo de termini, inter qui la ra-}
\l{\cel{7}porto aritmetikala (difero) o geometriala (quociento) esas}
\l{\cel{7}konstanta. - DEFIRS.}
\l{}
\l{\cel{4}\sou{projektar}. (\sou{trans.)}\cel{2}I. (\sou{geom.}) Trasar perpendikli de omna}
\l{\cel{7}punti di solido adsur plano, por figurar lu sur plano. -}
\l{\cel{7}II. (\sou{optiko})\cel{2}Reflektar sur+\sou{skrino,}per radii qui emanas de}
\l{\cel{7}lumo-foko. - DEFIRS.}
\l{}
\l{\cel{4}\sou{projektilo}. I. Irga korpo qua esas lansita per la impulso}
\l{\cel{7}da forco irgequala, por atingar ulu od ulo. - II. Korpo qua}
\l{\cel{7}esas lansita per lans-armo o per pafarmo. - DEFIS.}
\l{}
\l{\cel{4}\sou{projektoro}. (\sou{tekn.})\cel{2}Aparato uzata por reflektar, forme di un}
\l{\cel{7}o plura lumo-faski, di qui la intenseso lumala esas tre}
\l{\cel{7}granda, la lum-ondi de foko. - DEFIRS.}
\l{}
\l{\cel{4}\sou{projetar}. (\sou{trans.}) Pozar avane ideo kom realigenda. - DEFIRS.}
\l{}
\l{\cel{4}\sou{proklamar}. (\sou{trans.}) Anuncar publike e kelke-solene. - DEFIRS.}
\l{}
\l{\cel{4}\sou{proklitiko}. (\sou{gram.}) Vorto qua, pro indijar acento tonika, su}
\l{\cel{7}apogas sur la vorto qua sequas. - DEFIS.}
\l{}
\l{\cel{4}\sou{prokonsulo}.\cel{2}(en\cel{1}\sou{Roma, en la epoki antiqua}).\cel{2}Konsulo qua jus}
\l{\cel{7}ekiris de sua ofico-periodo ed a qua onu konfidis la komando}
\l{\cel{7}di armeo, di provinco. - DEFIRS.}
\l{}
\l{\cel{4}\sou{prokurar}. (\sou{trans.,}\cel{1}ad)\cel{2}Igar obtenar. - EFIS.}
\l{}
\l{\cel{4}\sou{prokuraco.}\cel{2}Darfo grantita legale da ulu a persono quan lu ko-}
\l{\cel{7}misas por agar ye la homo di la grantinto. - DEFIRS.}
}
